\documentclass[11pt]{article}
\usepackage{../homework}
\profname{Dr. Michael Martens}
\classcode{PHYS 313 Thermodynamics}
\assignnum{Assignment 04}
\duedate{07 October 2026}

\begin{document}
\begin{problem}{1}
A system in thermal contact with a reservoir at temperature \( \tau \) has states \( s \) with energies \( \varepsilon_s \).
The probability of finding the system in state \( s \) is
\begin{equation}
  \label{eq:p1-state-probability}
  P(\varepsilon_s) = \frac{e^{-\varepsilon_s/\tau}}{Z},
  \qquad \text{where} \qquad Z = \sum_s e^{-\varepsilon_s/\tau} = e^{-F/\tau}.
\end{equation}
\begin{enumerate}[label=(\alph*)]
  \item Show that the entropy of the system can be written as
  \begin{equation}
    \label{eq:p1-gibbs-entropy}
    \sigma = -\sum_s P(\varepsilon_s)\log P(\varepsilon_s).
  \end{equation}
  [Hint: Start with the expression for \( \log P(\varepsilon_s) \).]

  This expression for the entropy is useful when the probabilities of the system being in each of its states are known but the partition function is not.
        \begin{proof}
          As suggested in the hint, we first calculate
          \begin{equation}
          \label{eq:1}
          \begin{aligned}
            \log P(\varepsilon_s) &= \log( \frac{e^{-\varepsilon_s/\tau}}{Z} ) = \log(\exp(-\varepsilon_s/\tau)) - \log Z \\
            &= -\varepsilon_s/\tau - \sigma,
          \end{aligned}
        \end{equation}
        which implies
        \begin{equation}
        \label{eq:2}
        \sigma = -\varepsilon_s/\tau - \log(P(\varepsilon_{s})) = -\log(e^{\varepsilon_s/\tau} P(\varepsilon_{s}) )
        \end{equation}
        \end{proof}
\end{enumerate}
\end{problem}

The ideal gas (for Problems 2 and 3).
We did not cover the ideal gas in class.
A summary of its properties is given below.
For a monatomic gas of \( N \) particles of mass \( M \) in a volume \( V \) at temperature \( \tau \):
\begin{itemize}
  \item \( U = \frac{3}{2}N\tau \)
  \item \( pV = N\tau \)
  \item \( \sigma = N\bigl[\log(n_Q/n) + \frac{5}{2}\bigr] \), where \( n = N/V \) and \( n_Q = (M\tau/2\pi\hbar^2)^{3/2} \)
\end{itemize}
Also pay attention to the change in the entropy of the system for a reversible process.
If an isolated system undergoes a reversible process, its entropy \( \sigma_S \) remains constant.
If a system in contact with a thermal reservoir undergoes a reversible process, the total entropy \( \sigma_S + \sigma_R \) does not change, but the entropy of the system can change, with \( \Delta\sigma_S + \Delta\sigma_R = 0 \).

\begin{problem}{2}
A monatomic ideal gas is in thermal equilibrium with a reservoir at temperature \( \tau \).
The gas expands reversibly, its volume increasing by a small amount \( \delta V \), and it does work \( \delta W \) on its surroundings.
Find the change in the energy of the gas \( \delta U \), the heat \( \delta Q \) flowing into the gas, and the change in its free energy \( \delta F \).
Express each in terms of \( \delta W \).
\end{problem}

\begin{problem}{3}
The same gas is now thermally isolated.
It again expands reversibly, its volume increasing by \( \delta V \), and it does work \( \delta W \) on its surroundings.
Find \( \delta U \), \( \delta Q \), the temperature change \( \delta\tau \), and \( \delta F \).
Express each in terms of \( \delta W \).
\end{problem}

\begin{problem}{4}
Treat the Sun and the Earth as perfect blackbodies.
Use:
\begin{equation}
  \label{eq:p4-sun-earth-data}
  \begin{aligned}
    \text{Surface temperature of the Sun} \quad & T_{\odot} = 5800\,\mathrm{K} \\
    \text{Radius of the Sun} \quad & R_{\odot} = 6.96\times 10^8\,\mathrm{m} \\
    \text{Radius of the Earth} \quad & R_E = 6.37\times 10^6\,\mathrm{m} \\
    \text{Earth--Sun distance} \quad & d = 1.50\times 10^{11}\,\mathrm{m}
  \end{aligned}
\end{equation}
The Earth rotates fast enough that its surface temperature \( T_E \) may be taken as uniform.
Estimate \( T_E \).
\end{problem}

\begin{problem}{5}
A cavity of volume \( V \) is in thermal contact with a reservoir at temperature \( \tau \).
\begin{enumerate}[label=(\alph*)]
  \item Show that the number of thermal photons in the cavity is
  \begin{equation}
    \label{eq:p5-photon-number}
    N = \frac{2\zeta(3)}{\pi^2}V\left(\frac{\tau}{\hbar c}\right)^3
    \approx 0.244\,V\left(\frac{\tau}{\hbar c}\right)^3,
  \end{equation}
  where \( \zeta(3) = 1.202 \).
  You may use \( \displaystyle\int_0^{\infty}\frac{x^2\,dx}{e^x-1} = 2\zeta(3) \).

  \item The cosmic microwave background is blackbody radiation at \( T = 2.725\,\mathrm{K} \).
  Estimate the number of CMB photons per cubic centimeter.

  \item Using the Stefan--Boltzmann result \( U/V = \pi^2\tau^4/15\hbar^3c^3 \), find the average energy per photon in units of \( \tau \).
\end{enumerate}
\end{problem}

\begin{problem}{6}
A cubical cavity of volume \( V = L^3 \) contains thermal radiation.
The mode frequencies are
\begin{equation}
  \label{eq:p6-mode-frequencies}
  \omega_j = \frac{\pi c}{L}\sqrt{n_x^2+n_y^2+n_z^2},
\end{equation}
where \( j \) labels the modes.
A microstate of the radiation is specified by the number of photons \( s_j \) in each mode, and the total energy in the cavity is
\begin{equation}
  \label{eq:p6-cavity-energy}
  U = \sum_j s_j\hbar\omega_j.
\end{equation}
\begin{enumerate}[label=(\alph*)]
  \item The pressure is \( p = -(\partial U/\partial V)_{\sigma} \).
  If the volume is changed slowly enough, the occupation numbers \( s_j \) do not change, so the entropy stays constant.
  Use this to write the pressure of the photon gas as a sum over modes.

  \item Show that
  \begin{equation}
    \label{eq:p6-frequency-volume-derivative}
    \frac{d\omega_j}{dV} = -\frac{\omega_j}{3V}.
  \end{equation}

  \item Show that the pressure of thermal radiation is
  \begin{equation}
    \label{eq:p6-radiation-pressure}
    p = \frac{U}{3V}.
  \end{equation}
\end{enumerate}
\end{problem}

\begin{problem}{7}
Two large parallel black plates are held at temperatures \( T_1 \) (upper) and \( T_2 \) (lower), with \( T_1 > T_2 \).
The gap between them is evacuated, so energy crosses only by radiation.

See figure 1 in the original problem set (the unnumbered diagram in Problem 7).
The diagram labels the upper plate \( T_1 \), the intermediate shield \( T_s \) (part b), and the lower plate \( T_2 \).
\begin{enumerate}[label=(\alph*)]
  \item Find the net energy flux density \( J \) (power per unit area) from the upper plate to the lower one.

  \item A thin black plane---a heat shield---is placed between the plates, parallel to both.
  In the steady state, find the temperature \( T_s \) of the shield and the new net flux \( J_s \).

  \item A black surface at room temperature (\( 300\,\mathrm{K} \)) faces a black surface at liquid-helium temperature (\( 4\,\mathrm{K} \)).
  \begin{enumerate}[label=(\roman*)]
    \item Find the net energy flux density \( J \) (power per unit area) from the room-temperature surface to the liquid helium.

    \item A single heat shield is now placed between the room-temperature surface and the liquid helium.
    Find the equilibrium temperature of the shield and the new net energy flux density.
  \end{enumerate}
\end{enumerate}
\end{problem}
\end{document}
